\chapter{Sequences and Their Limits}

A capital of $1000$ euro earning $2\%$ interest per period does not grow in
one smooth sweep; it grows step by step: $1000$, $1020$, $1040.40,\dots$,
one number per period. Economic processes are full of such step-by-step
lists: iterated interest, instalment plans, prices adjusting period by
period towards an equilibrium, GDP figures year by year. The mathematical
object behind all of them is the \emph{sequence}, and the single most
important question about a sequence is what happens \emph{in the long run}:
does it settle down to a limit, or does it run off beyond every bound?

This chapter builds the toolkit. We state once, informally, what
convergence means, and then compute limits the practical way: with the
limit rules and a short list of standard limits that cover almost every
sequence met in economics. The star among them, $(1+x/n)^{n}\to\mathrm{e}^{x}$,
is exactly the step from period-by-period to \emph{continuous} compounding
and will carry the financial mathematics of a later chapter.

\section{Core Rules}

\begin{rulebox}{Sequences: explicit and recursive}
A \textbf{sequence} is a function that assigns to every index
$n\in\mathbb{N}$ a real number $a_{n}$; we write $(a_{n})_{n\geq 1}$ or
simply $a_{1},a_{2},a_{3},\dots$ It can be given
\begin{itemize}[leftmargin=1.6em,itemsep=0.3ex]
  \item \textbf{explicitly}, by a formula in $n$: e.g.\ $a_{n}=\dfrac{1}{n}$
        or $a_{n}=1000\cdot(1.02)^{n}$ (capital after $n$ periods at
        $2\%$ interest);
  \item \textbf{recursively}, by a start value and a rule for the next
        term: e.g.\ $a_{1}=1000$, $a_{n+1}=1.02\,a_{n}+100$ (interest plus
        a deposit of $100$ each period).
\end{itemize}
\end{rulebox}

\begin{rulebox}{Monotone and bounded sequences}
$(a_{n})$ is \textbf{increasing} if
$a_{n+1}\geq a_{n}$ for all $n$, and \textbf{decreasing} if
$a_{n+1}\leq a_{n}$ for all $n$ (with strict inequalities: strictly
increasing/decreasing). Either way the sequence is called
\textbf{monotone}.

$(a_{n})$ is \textbf{bounded above} if $a_{n}\leq M$ for some fixed $M$ and
all $n$, \textbf{bounded below} if $a_{n}\geq m$ for some fixed $m$, and
\textbf{bounded} if both hold.

\emph{Practical test for monotonicity:} compute the difference
$a_{n+1}-a_{n}$ and determine its sign; for sequences with positive terms
one may instead compare the ratio $a_{n+1}/a_{n}$ with $1$.
\end{rulebox}

\begin{rulebox}{Convergence and the limit}
$(a_{n})$ \textbf{converges} to the \textbf{limit} $a$ if the terms come arbitrarily close to $a$ and stay
there: for every tolerance $\varepsilon>0$, all but finitely many terms
satisfy $|a_{n}-a|<\varepsilon$. We write
\[
  \lim_{n\to\infty}a_{n}=a \qquad\text{or}\qquad a_{n}\to a .
\]
A sequence has \emph{at most one} limit. A sequence that does not converge
is \textbf{divergent}; if its terms eventually exceed every bound we say it
\emph{diverges to $\infty$} and write $a_{n}\to\infty$ (analogously
$-\infty$). In this book we never chase $\varepsilon$'s: limits are
computed with the rules and standard limits below.
\end{rulebox}

\begin{rulebox}{Limit rules}
Suppose $a_{n}\to a$ and $b_{n}\to b$. Then
\begin{itemize}[leftmargin=1.6em,itemsep=0.3ex]
  \item \textbf{Sum rule:} $a_{n}+b_{n}\to a+b$ \ (and $c\,a_{n}\to c\,a$
        for any constant $c$);
  \item \textbf{Product rule:} $a_{n}\,b_{n}\to a\,b$;
  \item \textbf{Quotient rule:} $\dfrac{a_{n}}{b_{n}}\to\dfrac{a}{b}$,
        provided $b\neq 0$;
  \item \textbf{Sandwich rule:} if $a_{n}\leq c_{n}\leq b_{n}$ from some
        index on and $a=b$, then $c_{n}\to a$ as well.
\end{itemize}
\emph{Working method for ratios of polynomials in $n$:} divide numerator
and denominator by the highest power of $n$ that occurs, then apply the
rules together with $1/n^{k}\to 0$.
\end{rulebox}

\begin{rulebox}{Standard limits}
\begin{itemize}[leftmargin=1.6em,itemsep=0.3ex]
  \item $\dfrac{1}{n^{k}}\to 0$ \ for every fixed $k>0$;
  \item $q^{n}\to 0$ \ if $|q|<1$; \quad $q^{n}\to\infty$ \ if $q>1$;
        \quad $q^{n}=1$ for $q=1$; \quad divergent for $q\leq -1$;
  \item $\sqrt[n]{n}\to 1$ \ and \ $\sqrt[n]{c}\to 1$ for every constant
        $c>0$;
  \item $\Bigl(1+\dfrac{x}{n}\Bigr)^{n}\to\mathrm{e}^{x}$ \ for every
        $x\in\R$; in particular
        $\Bigl(1+\dfrac{1}{n}\Bigr)^{n}\to\mathrm{e}\approx 2.71828$.
\end{itemize}
The last limit is the bridge to \emph{continuous compounding}: splitting a
yearly interest rate $x$ into $n$ compounding steps turns capital $K$ into
$K\,(1+x/n)^{n}$, and as the steps become infinitely fine this tends to
$K\,\mathrm{e}^{x}$, the formula behind continuous interest in financial
mathematics.
\end{rulebox}

\begin{rulebox}{Monotone convergence principle}
\textbf{Bounded $+$ monotone $\Rightarrow$ convergent.} An increasing
sequence that is bounded above converges; a decreasing sequence that is
bounded below converges. This is the workhorse for recursive sequences:
first establish monotonicity and boundedness, then find the candidate
limit $a$ from the \emph{fixed-point equation} obtained by letting
$n\to\infty$ in the recursion (e.g.\ $a=1.02\,a+100$ for
$a_{n+1}=1.02\,a_{n}+100$).
\end{rulebox}

\begin{rulebox}{Growth hierarchy}
For every $k>0$ and every $q>1$,
\[
  \ln n \;\ll\; n^{k} \;\ll\; q^{n},
  \qquad\text{meaning}\qquad
  \frac{\ln n}{n^{k}}\to 0
  \quad\text{and}\quad
  \frac{n^{k}}{q^{n}}\to 0 .
\]
Every power beats every logarithm, and every exponential with base $q>1$
beats every power. A useful consequence: $n^{k}q^{n}\to 0$ whenever
$|q|<1$, geometric decay swallows polynomial growth.
\end{rulebox}

\section{Practice Problems}

Compute the limit of each sequence or show that it diverges; where
indicated, also decide monotonicity and boundedness.

\begin{problems}
  \item $a_{n}=\dfrac{3n+5}{n+2}$
  \item $a_{n}=\dfrac{2n^{2}-n+4}{5n^{2}+3}$
  \item $a_{n}=\dfrac{n+1}{n^{2}+1}$
  \item $a_{n}=\dfrac{n^{3}-2n}{4n^{2}+n}$
  \item $a_{n}=\sqrt{n^{2}+n}-n$
  \item $a_{n}=\dfrac{\sqrt{4n^{2}+1}}{3n-2}$
  \item $a_{n}=(0.9)^{n}$ \ and \ $b_{n}=(1.1)^{n}$ \quad (limit or
        divergence for each)
  \item $a_{n}=\dfrac{2^{n}+5^{n}}{5^{n}}$
  \item $a_{n}=\dfrac{3^{n}-2^{n}}{3^{n}+2^{n}}$
  \item $a_{n}=\sqrt[n]{5n}$
  \item $a_{n}=\Bigl(1+\dfrac{3}{n}\Bigr)^{n}$
  \item $a_{n}=\Bigl(1-\dfrac{2}{n}\Bigr)^{n}$
  \item $a_{n}=\Bigl(1+\dfrac{1}{2n}\Bigr)^{n}$
  \item $a_{n}=\Bigl(\dfrac{n+4}{n}\Bigr)^{n}$
  \item $a_{n}=\dfrac{n}{n+1}$ \quad (monotone? bounded? convergent?,
        justify all three)
  \item $b_{n}=(-1)^{n}\,\dfrac{n}{n+1}$ \quad (monotone? bounded?
        convergent?)
  \item Savings account: $a_{1}=1000$ and $a_{n+1}=1.02\,a_{n}+100$
        (interest of $2\%$ per period plus a deposit of $100$). Compute
        $a_{2}$ and $a_{3}$, derive a closed formula for $a_{n}$, and
        describe the behaviour as $n\to\infty$.
  \item $a_{n}=\dfrac{2n+(-1)^{n}}{n}$ \quad (use the sandwich rule)
  \item $a_{n}=\dfrac{n^{2}}{2^{n}}$
  \item Classify as convergent or divergent, giving the limit where it
        exists:
        \[
          \text{(a)}\; a_{n}=(-1)^{n},\qquad
          \text{(b)}\; b_{n}=5-\frac{3}{n},\qquad
          \text{(c)}\; c_{n}=\ln n,\qquad
          \text{(d)}\; d_{n}=(-0.5)^{n}.
        \]
\end{problems}

\section{Solutions}

\begin{solutions}
  \item Divide numerator and denominator by $n$:
        \[
          a_{n}=\frac{3+\dfrac{5}{n}}{1+\dfrac{2}{n}}
          \;\longrightarrow\;\frac{3+0}{1+0}=3 ,
        \]
        using $5/n\to 0$, $2/n\to 0$ and the quotient rule. So
        $\lim_{n\to\infty}a_{n}=3$.

  \item Divide by the highest power $n^{2}$:
        \[
          a_{n}=\frac{2-\dfrac{1}{n}+\dfrac{4}{n^{2}}}{5+\dfrac{3}{n^{2}}}
          \;\longrightarrow\;\frac{2-0+0}{5+0}=\frac{2}{5} .
        \]

  \item Divide by $n^{2}$ (the highest power in the denominator):
        \[
          a_{n}=\frac{\dfrac{1}{n}+\dfrac{1}{n^{2}}}{1+\dfrac{1}{n^{2}}}
          \;\longrightarrow\;\frac{0+0}{1+0}=0 .
        \]
        Degree of the numerator below degree of the denominator: limit $0$.

  \item Divide by $n^{2}$:
        \[
          a_{n}=\frac{n-\dfrac{2}{n}}{4+\dfrac{1}{n}} .
        \]
        The denominator tends to $4$, while the numerator $n-2/n\to\infty$.
        Hence $a_{n}\to\infty$: the sequence is \textbf{divergent}
        (for large $n$, $a_{n}\approx n/4$).

  \item Multiply by the conjugate:
        \[
          a_{n}=\bigl(\sqrt{n^{2}+n}-n\bigr)
                \cdot\frac{\sqrt{n^{2}+n}+n}{\sqrt{n^{2}+n}+n}
               =\frac{(n^{2}+n)-n^{2}}{\sqrt{n^{2}+n}+n}
               =\frac{n}{\sqrt{n^{2}+n}+n} .
        \]
        Now divide numerator and denominator by $n$:
        \[
          a_{n}=\frac{1}{\sqrt{1+\dfrac{1}{n}}+1}
          \;\longrightarrow\;\frac{1}{\sqrt{1}+1}=\frac{1}{2} .
        \]

  \item Divide numerator and denominator by $n$ (inside the root: by
        $n^{2}$):
        \[
          a_{n}=\frac{\sqrt{4+\dfrac{1}{n^{2}}}}{3-\dfrac{2}{n}}
          \;\longrightarrow\;\frac{\sqrt{4}}{3}=\frac{2}{3} .
        \]

  \item For $a_{n}=(0.9)^{n}$: here $q=0.9$ with $|q|<1$, so
        $a_{n}\to 0$ (convergent). For $b_{n}=(1.1)^{n}$: here $q=1.1>1$,
        so $b_{n}\to\infty$ (divergent), capital growing at $10\%$ per
        period exceeds every bound.

  \item Split the fraction:
        \[
          a_{n}=\frac{2^{n}}{5^{n}}+\frac{5^{n}}{5^{n}}
               =\Bigl(\frac{2}{5}\Bigr)^{n}+1
          \;\longrightarrow\;0+1=1 ,
        \]
        since $|2/5|<1$.

  \item Divide numerator and denominator by the dominant term $3^{n}$:
        \[
          a_{n}=\frac{1-\bigl(\tfrac{2}{3}\bigr)^{n}}
                     {1+\bigl(\tfrac{2}{3}\bigr)^{n}}
          \;\longrightarrow\;\frac{1-0}{1+0}=1 ,
        \]
        because $(2/3)^{n}\to 0$ ($|q|<1$).

  \item Factor the root:
        \[
          a_{n}=\sqrt[n]{5n}=\sqrt[n]{5}\cdot\sqrt[n]{n}
          \;\longrightarrow\;1\cdot 1=1 ,
        \]
        by the standard limits $\sqrt[n]{c}\to 1$ ($c>0$) and
        $\sqrt[n]{n}\to 1$, combined with the product rule.

  \item This is the standard limit with $x=3$:
        \[
          \lim_{n\to\infty}\Bigl(1+\frac{3}{n}\Bigr)^{n}=\mathrm{e}^{3} .
        \]

  \item Standard limit with $x=-2$:
        \[
          \lim_{n\to\infty}\Bigl(1-\frac{2}{n}\Bigr)^{n}
          =\mathrm{e}^{-2}=\frac{1}{\mathrm{e}^{2}} .
        \]

  \item Rewrite the inner fraction so the pattern $x/n$ appears:
        $\dfrac{1}{2n}=\dfrac{1/2}{n}$, so
        \[
          a_{n}=\Bigl(1+\frac{1/2}{n}\Bigr)^{n}
          \;\longrightarrow\;\mathrm{e}^{1/2}=\sqrt{\mathrm{e}} .
        \]

  \item Simplify the base first:
        \[
          \frac{n+4}{n}=1+\frac{4}{n}
          \qquad\Longrightarrow\qquad
          a_{n}=\Bigl(1+\frac{4}{n}\Bigr)^{n}\to\mathrm{e}^{4} .
        \]

  \item \textbf{Monotone:} compute the difference
        \[
          a_{n+1}-a_{n}=\frac{n+1}{n+2}-\frac{n}{n+1}
          =\frac{(n+1)^{2}-n(n+2)}{(n+1)(n+2)}
          =\frac{n^{2}+2n+1-n^{2}-2n}{(n+1)(n+2)}
          =\frac{1}{(n+1)(n+2)}>0 ,
        \]
        so $(a_{n})$ is strictly increasing.
        \textbf{Bounded:} $0<a_{n}<1$ for all $n$, since $n<n+1$.
        \textbf{Convergent:} yes, by the monotone convergence principle
        (increasing and bounded above), and directly:
        \[
          a_{n}=\frac{1}{1+\dfrac{1}{n}}\;\longrightarrow\;\frac{1}{1+0}=1 .
        \]

  \item First terms: $b_{1}=-\tfrac{1}{2}$, $b_{2}=\tfrac{2}{3}$,
        $b_{3}=-\tfrac{3}{4}$, $b_{4}=\tfrac{4}{5}$. The signs alternate,
        so the sequence is \textbf{not monotone}. It is \textbf{bounded},
        because $|b_{n}|=\dfrac{n}{n+1}<1$ for all $n$. It is
        \textbf{divergent}: the even-indexed terms
        $b_{2m}=\tfrac{2m}{2m+1}\to 1$ while the odd-indexed terms tend to
        $-1$; a convergent sequence cannot approach two different values
        (uniqueness of the limit). Boundedness alone does not force
        convergence.

  \item First terms:
        \[
          a_{2}=1.02\cdot 1000+100=1020+100=1120 ,\qquad
          a_{3}=1.02\cdot 1120+100=1142.40+100=1242.40 .
        \]
        \textbf{Fixed point:} a constant solution $a^{*}$ would satisfy
        $a^{*}=1.02\,a^{*}+100$, i.e.\ $-0.02\,a^{*}=100$, so
        $a^{*}=-5000$. \textbf{Closed formula:} set
        $d_{n}=a_{n}-a^{*}=a_{n}+5000$. Then
        \[
          d_{n+1}=a_{n+1}+5000=1.02\,a_{n}+100+5000
                 =1.02\,(a_{n}+5000)=1.02\,d_{n} ,
        \]
        a geometric sequence with $d_{1}=1000+5000=6000$. Hence
        $d_{n}=6000\cdot(1.02)^{n-1}$ and
        \[
          a_{n}=6000\cdot(1.02)^{n-1}-5000 .
        \]
        Check: $a_{3}=6000\cdot 1.0404-5000=6242.40-5000=1242.40$,
        matches. \textbf{Behaviour:} $(a_{n})$ is strictly increasing
        ($a_{n+1}-a_{n}=0.02\,a_{n}+100>0$ since all terms are positive)
        and, because $1.02>1$ gives $(1.02)^{n-1}\to\infty$, the sequence
        \textbf{diverges to $\infty$}: interest plus steady deposits grow
        beyond every bound.

  \item Split: $a_{n}=2+\dfrac{(-1)^{n}}{n}$. Since
        $-1\leq(-1)^{n}\leq 1$,
        \[
          2-\frac{1}{n}\;\leq\;a_{n}\;\leq\;2+\frac{1}{n} .
        \]
        Both bounds tend to $2$, so by the sandwich rule
        $\lim_{n\to\infty}a_{n}=2$.

  \item Growth hierarchy: $n^{k}\ll q^{n}$ for $q>1$, here with $k=2$ and
        $q=2$, so
        \[
          \lim_{n\to\infty}\frac{n^{2}}{2^{n}}=0 .
        \]
        The decay is fast: for $n=10$ the term is
        $\tfrac{100}{1024}\approx 0.098$, for $n=20$ already
        $\tfrac{400}{1048576}\approx 0.0004$. Convergent with limit $0$.

  \item (a) $(-1)^{n}$ alternates between $-1$ and $1$ forever: bounded
        but \textbf{divergent} (no single limit). \
        (b) $3/n\to 0$, so $b_{n}\to 5-0=5$: \textbf{convergent} with
        limit $5$. \
        (c) $\ln n$ grows without bound (slowly, but beyond every $M$):
        \textbf{divergent} to $\infty$. \
        (d) $q=-0.5$ with $|q|<1$, so $d_{n}\to 0$: \textbf{convergent}
        with limit $0$.
\end{solutions}

\section{Additional Problems}

\begin{problems}
  \item $a_{n}=\dfrac{7n-4}{2n+9}$
  \item $a_{n}=\dfrac{n^{2}+6n}{3n^{2}-n+1}$
  \item $a_{n}=\dfrac{\sqrt{9n^{2}+n}}{n+1}$
  \item $a_{n}=\sqrt{n+1}-\sqrt{n}$
  \item $a_{n}=\dfrac{4^{n}-3^{n}}{4^{n}}$
  \item $a_{n}=\Bigl(1+\dfrac{5}{n}\Bigr)^{n}$
  \item $a_{n}=\Bigl(1-\dfrac{1}{n}\Bigr)^{2n}$
  \item $a_{n}=3+\dfrac{2}{n}$ \quad (monotone? bounded? convergent?,
        justify all three)
  \item Price adjustment: $p_{1}=200$ and $p_{n+1}=0.5\,p_{n}+60$ (each
        period the price moves halfway towards a demand-driven level plus
        a constant mark-up). Compute $p_{2}$ and $p_{3}$, derive a closed
        formula for $p_{n}$, and describe the behaviour as $n\to\infty$.
  \item Classify as convergent or divergent, giving the limit where it
        exists:
        \[
          \text{(a)}\; a_{n}=\frac{n}{2^{n}},\qquad
          \text{(b)}\; b_{n}=(1.05)^{n},\qquad
          \text{(c)}\; c_{n}=(-1)^{n}\,n,\qquad
          \text{(d)}\; d_{n}=\frac{2n+3}{n+1}.
        \]
\end{problems}

\section{Additional Solutions}

\begin{solutions}
  \item Divide numerator and denominator by $n$:
        \[
          a_{n}=\frac{7-\dfrac{4}{n}}{2+\dfrac{9}{n}}
          \;\longrightarrow\;\frac{7-0}{2+0}=\frac{7}{2} .
        \]

  \item Divide by $n^{2}$:
        \[
          a_{n}=\frac{1+\dfrac{6}{n}}{3-\dfrac{1}{n}+\dfrac{1}{n^{2}}}
          \;\longrightarrow\;\frac{1+0}{3-0+0}=\frac{1}{3} .
        \]

  \item Divide numerator and denominator by $n$; under the root this means
        dividing by $n^{2}$:
        \[
          a_{n}=\frac{\sqrt{9+\dfrac{1}{n}}}{1+\dfrac{1}{n}}
          \;\longrightarrow\;\frac{\sqrt{9}}{1}=3 .
        \]

  \item Multiply by the conjugate:
        \[
          a_{n}=\bigl(\sqrt{n+1}-\sqrt{n}\bigr)
                \cdot\frac{\sqrt{n+1}+\sqrt{n}}{\sqrt{n+1}+\sqrt{n}}
               =\frac{(n+1)-n}{\sqrt{n+1}+\sqrt{n}}
               =\frac{1}{\sqrt{n+1}+\sqrt{n}} .
        \]
        The denominator tends to $\infty$, hence $a_{n}\to 0$.

  \item Split the fraction:
        \[
          a_{n}=1-\Bigl(\frac{3}{4}\Bigr)^{n}
          \;\longrightarrow\;1-0=1 ,
        \]
        since $|3/4|<1$.

  \item Standard limit with $x=5$:
        \[
          \lim_{n\to\infty}\Bigl(1+\frac{5}{n}\Bigr)^{n}=\mathrm{e}^{5} .
        \]

  \item Write the power as a square of the standard form:
        \[
          a_{n}=\Bigl[\Bigl(1-\frac{1}{n}\Bigr)^{n}\Bigr]^{2}
          \;\longrightarrow\;\bigl(\mathrm{e}^{-1}\bigr)^{2}
          =\mathrm{e}^{-2}=\frac{1}{\mathrm{e}^{2}} ,
        \]
        using $(1-1/n)^{n}\to\mathrm{e}^{-1}$ and the product rule.

  \item \textbf{Monotone:}
        \[
          a_{n+1}-a_{n}=\frac{2}{n+1}-\frac{2}{n}
          =\frac{2n-2(n+1)}{n(n+1)}=\frac{-2}{n(n+1)}<0 ,
        \]
        so $(a_{n})$ is strictly decreasing.
        \textbf{Bounded:} $3<a_{n}\leq a_{1}=5$ for all $n$.
        \textbf{Convergent:} yes: decreasing and bounded below, so the
        monotone convergence principle applies; since $2/n\to 0$, the
        limit is $3$.

  \item First terms:
        \[
          p_{2}=0.5\cdot 200+60=100+60=160 ,\qquad
          p_{3}=0.5\cdot 160+60=80+60=140 .
        \]
        \textbf{Fixed point:} $p^{*}=0.5\,p^{*}+60$ gives
        $0.5\,p^{*}=60$, so $p^{*}=120$. \textbf{Closed formula:} set
        $d_{n}=p_{n}-120$. Then
        \[
          d_{n+1}=p_{n+1}-120=0.5\,p_{n}+60-120=0.5\,(p_{n}-120)
                 =0.5\,d_{n} ,
        \]
        with $d_{1}=200-120=80$, so $d_{n}=80\cdot(0.5)^{n-1}$ and
        \[
          p_{n}=120+80\cdot(0.5)^{n-1} .
        \]
        Check: $p_{3}=120+80\cdot 0.25=120+20=140$, matches.
        \textbf{Behaviour:} $(0.5)^{n-1}\to 0$ because $|0.5|<1$, so
        $p_{n}\to 120$: the price falls monotonically to the equilibrium
        $120$. Contrast with the savings account of Problem~17 in the main
        list: a coefficient below $1$ pulls the sequence \emph{towards}
        the fixed point, a coefficient above $1$ (there $1.02$) pushes it
        \emph{away}.

  \item (a) Growth hierarchy ($n^{k}\ll 2^{n}$ with $k=1$):
        $a_{n}\to 0$, \textbf{convergent}. \
        (b) $q=1.05>1$, so $b_{n}\to\infty$: \textbf{divergent}. \
        (c) $|c_{n}|=n\to\infty$ with alternating signs: unbounded, hence
        \textbf{divergent} (no limit, not even $\pm\infty$). \
        (d) Divide by $n$:
        $d_{n}=\dfrac{2+3/n}{1+1/n}\to\dfrac{2}{1}=2$:
        \textbf{convergent} with limit $2$.
\end{solutions}
